\documentstyle[12pt,a4,german,epsf,twoside]{article}

\newcommand{\nl}{{\newline}}

\title{Diplomarbeit - "Ubersicht}
\author{Patrick Stein - Matrikelnummer 814212}
\begin{document}
\maketitle

\section{Thema: Verbesserung von Suchmaschinen im World\/Wide\/Web}

Die momentane Informationsrechere im im World\/Wide\/Web (WWW) ist zum Gro"steil recht m"uhselig. Es wird dem, Informationen suchenden, Benutzer recht schwer gemacht, genau die gesuchten Informationen zu finden. Die Rechere im {\it WWW} wird zum jetzigen Zeitpunkt nur von einfachen Indizes (Suchmaschinen) erleichtert. Diese bieten meist nur eine Schlagwortrechere an, so da"s die Antworten, die diese Indizes liefern, meist enorm viele Treffer liefern. Da jedoch die Treffer nur entfernt mit dem gesuchten Begriff zu tun haben, sollte man versuchen, die Treffermengen einzugrenzen.\nl
In der Diplomarbeit wird versucht werden, Methoden zu finden und zu implementieren, die gute Aussichten haben, das Problem zu l"osen:
\begin{description}

	\item[Sprachenerkennung] Wenn es m"oglich ist, die Sprache(n) zu erkennen, in der ein Dokument geschrieben ist, l"a"st sich f"ur den Informationssuchenden einfach einstellen, welche Dokumente f"ur Ihn unn"utz, da unverst"andlich, sind.\nl
	Beispiel: Die Suche nach {\it Film} liefert bei den jetzigen Suchindizes sehr viele amerikanische und englische Datenquellen. K"onnte man einfach sagen {\it Zeige mir alle in Deutsch geschriebenen Seiten in denen das Wort {\it "`Film"'} vorkommt} w"are die Suche schon wesentlich effektiver (f"ur den Fall, man will wirklich nur nach Seiten in Deutscher Sprache suchen).

	\item[Zusammenh"ange] Die Qualit"at der Suche ist im {\it WWW} eigentlich nicht besser, als die einer gut ausgestatteten Bibliothek. Sucht man zu Beispiel etwas "uber {\it Mozart}, so offenbart die heute g"angige Bibliothek einem eine F"ulle von B"uchern in denen {\it Mozart} abgehandelt wird oder man erf"ahrt aus einem Lexikon die wichtigsten Daten.\nl
	 Die Suche im {\it WWW} geht schon etwas weiter: Die Daten "uber {\it Mozart} tauchen auch dort auf, jedoch erh"alt der Benutzer weit mehr Querverweise als in Printmedien.\nl
	 Noch besser jedoch w"are, die gefundenen Daten in {\it Clusteren} aufzubereiten. So k"onnte man sich vorstellen, das bei einer Suche nach {\it Mozart} die Ergebnismenge {\it Mozart, Beethoven, klassische Musik, Die kleine Nachtmusik, "Osterreich,...} liefert.

	\item[{\it Visualisierung}] Die gefundenen {\it Cluster} k"onnten am Bildschirm in einer Art Graph dargestellt werden in dem man dann sich durch Themengebiete sehr viel effektiver arbeiten w"urde.
\end{description}

\subsection{Pr"ufer}
Prof. Dr. Kirch-Prinz		\nl
Fachbereich Informatik		\nl
Fachhochschule M"unchen		\nl
Lothstra"se 34				\nl
80 335 M"unchen				\nl
Tel. (089) 2165 - 1639		\nl
\nl
Prof. Dr. G"unthner			\nl
Centrum f"ur Informations- und Sprachverarbeitung	\nl
Ludwig-Maximiliansuniversit"at M"unchen	\nl
"Ottingenstra"se 67			\nl
80 538 M"unchen				\nl
Tel. (089) 2178 - 2721		\nl

\subsection{Diplomand}
Patrick Stein	\nl
Alte M"unchener Stra"se 50	\nl
85 774 M"unchen	\nl
Matrikelnummer: 814212	\nl
Tel. (089) 2178 - 2715 ( im CIS - Uni M"unchen )	\nl
Tel. (089) 950 57 34 ( zu Hause )	\nl

\end{document}







